% translation-source: docs/user/latex/chapters/06-courbes.tex; sha256: c9120e88bac82ef65f76b3799e02874eb224ebbb0f5c048893e981b5bf2a188e
\chapter{Tone curve}
\section{Choose a channel}
Select luminance or the red, green or blue channel using the module buttons.
The diagonal represents the reference mapping. Editing a channel can introduce
a color cast even if its curve resembles the luminance curve.

\section{Manipulate points}
Double-click the graph to add a handle, drag it to change the curve, and
right-click an interior handle to remove it. Endpoints do not behave like
interior points. Avoid accumulating closely spaced points: a simpler curve
is easier to control.

Use reset for the current curve or all curves as appropriate. Double-clicking
this graph \emph{adds a point}: it is not the sliders' reset gesture.
No dedicated keyboard shortcut for adding/removing points is declared.

\section{Monochrome background}
The gray silhouette shows the current preview's luminance and stays monochrome
when an RGB curve is selected. It is a post-render reference, not the exact
histogram of the linear signal entering the curves.

\paragraph{Example.}
For gentle contrast, start with two widely spaced interior points, lower
the dark region slightly and raise the bright region slightly. Compare using
\touche{Ctrl + Shift + Y}; watch transitions rather than seeking a perfect
geometric shape.
